%%%%%%%%%%%%%%%%%%%%%%%%%%%%%%%%%%%%%%%%%
%----------------------------------------------------------------------------------------
%	PACKAGES AND OTHER DOCUMENT CONFIGURATIONS
%----------------------------------------------------------------------------------------

\documentclass[a4paper,12pt]{article} % добавить leqno в [] для нумерации слева

%%% Работа с русским языком
\usepackage{cmap}					% поиск в PDF
\usepackage{mathtext} 				% русские буквы в формулах
\usepackage[T2A]{fontenc}			% кодировка
\usepackage[utf8]{inputenc}			% кодировка исходного текста
\usepackage[english,russian]{babel}	% локализация и переносы
\usepackage{graphicx}

%%% Дополнительная работа с математикой
\usepackage{amsmath,amsfonts,amssymb,amsthm,mathtools} % AMS
\usepackage{icomma} % "Умная" запятая: $0,2$ --- число, $0, 2$ --- перечисление

%% Номера формул
%\mathtoolsset{showonlyrefs=true} % Показывать номера только у тех формул, на которые есть \eqref{} в тексте.

%% Шрифты
\usepackage{euscript}	 % Шрифт Евклид
\usepackage{mathrsfs} % Красивый матшрифт

%% Перенос знаков в формулах (по Львовскому)
\newcommand*{\hm}[1]{#1\nobreak\discretionary{}
{\hbox{$\mathsurround=0pt #1$}}{}}

%%% Работа с картинками
\usepackage{graphicx}  % Для вставки рисунков
\graphicspath{{images/}{images2/}}  % папки с картинками
\setlength\fboxsep{3pt} % Отступ рамки \fbox{} от рисунка
\setlength\fboxrule{1pt} % Толщина линий рамки \fbox{}
\usepackage{wrapfig} % Обтекание рисунков и таблиц текстом

%%% Работа с таблицами
\usepackage{array,tabularx,tabulary,booktabs} % Дополнительная работа с таблицами
\usepackage{longtable}  % Длинные таблицы
\usepackage{multirow} % Слияние строк в таблице

% Переменные
\newcommand{\MAIITitle}{Московский Авиационный Институт \\ (Национальный исследовательский университет)}
\newcommand{\MAIInstitute}{8}
\newcommand{\MAIInstituteName}{Информационные технологии и прикладная математика}
\newcommand{\MAIDepartment}{813}
\newcommand{\MAIDepartmentName}{Компьютерная математика}

\newcommand{\MAIStudentName}{Васильев Дмитрий Олегович}
\newcommand{\MAITeacherName}{Родников Александр Владимирович}
\newcommand{\MAIGroup}{М8О-310Б-18}
\newcommand{\MAIDate}{1.01.12}


\begin{document}

%----------------------------------------------------------------------------------------
%	TITLE PAGE
%----------------------------------------------------------------------------------------

\begin{titlepage}
	\newcommand{\HRule}{\rule{\linewidth}{0.5mm}} % Defines a new command for horizontal lines, change thickness here
	
	\center
	
	%------------------------------------------------
	%	Заголовок
	%------------------------------------------------
	
	\textsc{\large \MAIITitle}\\[1.25cm] % Main heading such as the name of your university/college
	
	%------------------------------------------------
	%	Логотип института
	%------------------------------------------------
	
	\includegraphics[width=2.5cm]{mai-logo.eps}
	
	%----------------------------------------------------------------------------------------
	
	\textsc{\Large Институт №\MAIInstitute \\ «\MAIInstituteName»}\\[0.5cm] % Название института
	
	\textsc{\large Кафедра \MAIDepartment \\ «\MAIDepartmentName»}\\[0.5cm] % Название кафедры
	
	%------------------------------------------------
	%	Название
	%------------------------------------------------
	
	\HRule\\[0.4cm]
	
	{\huge\bfseries Применение фундаментальных алгоритмов для анализа данных}\\[0.5cm]
	
	\HRule\\[1.5cm]
	
	%------------------------------------------------
	%	Студент, группа, преподаватель, оценка, дата.
	%------------------------------------------------
	
	\vfill\vfill
	
	\begin{minipage}{1\textwidth}
		\begin{flushright}
			\begin{tabular}{|l|r|}
				\hline 
				Студент & \MAIStudentName \\ 
				\hline 
				Группа & \MAIGroup \\ 
				\hline 
				Преподаватель & \MAITeacherName \\ 
				\hline 
				Оценка & \\ 
				\hline 
				Дата & \MAIDate \\ 
				\hline 
			\end{tabular} 
		\end{flushright}
	\end{minipage}
	
	
	%------------------------------------------------
	%	Date
	%------------------------------------------------
	
	\vfill\vfill\vfill
	
	{\large Москва \the\year{}} 
	
\end{titlepage}

%----------------------------------------------------------------------------------------

\end{document}